% ch13_categories_as_algebra.tex — Volume II Chapter 1

\chapter{Categories as Algebraic Structures}

The first volume introduced categories from the outside: objects, morphisms,
composition, identity. This chapter looks at them from the inside, as algebraic
structures in the sense of abstract algebra. The key insight is that a category
is a \emph{monoid with multiple objects} — and that monoids, groups, preorders,
and groupoids are each a special type of category.

This algebraic viewpoint sharpens the intuition for what is ``categorically
natural'' versus what is a coincidence of a specific example.

% ─────────────────────────────────────────────────────────────────────────────
\begin{nameorigin}
\term{Monoid} (Greek \emph{monos} ``single, one'' + \emph{-oid}
``having the form of'') is a set with an associative multiplication and
a single (= ``mon-'') identity element. The term is 20th-century
algebraic terminology, standard by the 1930s. \term{Groupoid} (Heinrich
Brandt 1927) is the analogue with all elements invertible — and crucially,
where the multiplication is \emph{partial} (not every pair multiplies).
Brandt's insight: a groupoid is to a group what a category is to a
monoid. The slogan ``a category is a many-object monoid'' captures this
exactly.
\end{nameorigin}

\section{Categories as Generalized Monoids}

A monoid is a set $M$ with an associative binary operation and a unit element.
A category has objects, morphisms, associative composition, and identity morphisms.
The difference: in a monoid, any two elements can be multiplied. In a category,
two morphisms $f : a \to b$ and $g : c \to d$ can only be composed when $b = c$.
Composition is \emph{partial}.

\begin{dialog}
``A category is just a monoid where the multiplication might not be defined?''

Exactly. The technical term is \term{monoidoid} — rarely used — but the idea is
precise: a category is to a monoid what a groupoid is to a group.
\end{dialog}

The correspondence is made exact by passing to one-object categories.

\begin{theorem}
There is a bijection between monoids and categories with exactly one object.
\end{theorem}

\begin{prooflog}
\proofstep{Let $\mathcal{C}$ be a category with one object, say $\star$.}
           {$\mathcal{C}$ has exactly the morphisms $\mathcal{C}(\star, \star)$.}
\proofstep{Composition $\circ : \mathcal{C}(\star,\star) \times \mathcal{C}(\star,\star) \to \mathcal{C}(\star,\star)$
           is total, since all morphisms share domain and codomain $\star$.}
           {In a one-object category, every morphism has both domain and codomain equal to $\star$.}
\proofstep{Composition is associative (category axiom); $\operatorname{id}_\star$ is
           a two-sided unit.}
           {This is precisely the monoid axioms for $(\mathcal{C}(\star,\star), \circ, \operatorname{id}_\star)$.}
\proofstep{Conversely, given a monoid $(M, \cdot, e)$, define a category $\mathbf{B}M$
           with one object $\star$ and morphisms $\mathbf{B}M(\star,\star) = M$.}
           {$\mathbf{B}$ stands for ``delooping'' — a standard construction.}
\proofstep{Composition in $\mathbf{B}M$ is multiplication in $M$; identity is $e$.
           Associativity and unit laws hold because $M$ is a monoid.}
           {The two constructions are inverses: $\mathbf{B}M$ round-trips a monoid
            back to itself.}
\end{prooflog}

\begin{howtosay}
$\mathbf{B}M$ is read ``the delooping of $M$'' or ``$M$ viewed as a one-object category.''
The ``B'' comes from the classifying space construction in topology.
\end{howtosay}

The same pattern gives a bijection between groups and one-object groupoids (categories
where every morphism is an isomorphism).

% ─────────────────────────────────────────────────────────────────────────────
\section{Preorders, Groupoids, and the Algebraic Taxonomy}

A \term{preorder} $(P, \leq)$ gives a category: objects are elements, and there
is exactly one morphism $a \to b$ if $a \leq b$, no morphism otherwise. The
category axioms become: $a \leq a$ (reflexivity = identity) and $a \leq b$,
$b \leq c$ implies $a \leq c$ (transitivity = composition). A preorder is
a category where each hom-set has at most one element.

A \term{groupoid} is a category where every morphism is an isomorphism.
Groups are one-object groupoids. Equivalence relations, viewed as preorders,
are groupoids that are also preorders: since $a \sim b \Leftrightarrow b \sim a$,
every morphism $a \to b$ has an inverse $b \to a$.

\begin{center}
\renewcommand{\arraystretch}{1.4}
\begin{tabular}{lll}
\hline
\textbf{Algebraic structure} & \textbf{As a category} & \textbf{Condition} \\
\hline
Monoid & 1-object category & --- \\
Group & 1-object groupoid & every morphism invertible \\
Preorder & thin category & at most 1 morphism per hom-set \\
Partial order & thin category with no cycles & thin + no $a \neq b$ with $a \leq b \leq a$ \\
Groupoid & category & every morphism invertible \\
Equivalence relation & thin groupoid & thin + invertible \\
\hline
\end{tabular}
\end{center}

These are not just analogies — they are literal special cases. The category
$\mathbf{Mon}$ of monoids is a full subcategory of $\mathbf{Cat}$ via $\mathbf{B}$.
The category $\mathbf{Grp}$ embeds into $\mathbf{Gpd}$ (groupoids). The category
$\mathbf{Pos}$ of partial orders embeds into $\mathbf{Cat}$.

% ─────────────────────────────────────────────────────────────────────────────
\section{Ab-Enriched Categories and the Additive Setting}

A category is \term{$\mathbf{Ab}$-enriched} when each hom-set $\mathcal{C}(a,b)$
is an abelian group and composition is bilinear:
\[
  (f + f') \circ g = f \circ g + f' \circ g,
  \qquad
  f \circ (g + g') = f \circ g + f \circ g'.
\]
Such categories are called \term{preadditive}. The zero morphism $0_{a,b} : a \to b$
is the zero element of the abelian group $\mathcal{C}(a,b)$.

When a preadditive category also has a zero object ($0$ with $\mathcal{C}(0,a)$ and
$\mathcal{C}(a,0)$ each having one element for all $a$), and all finite products
and coproducts, it is called \term{additive}. In additive categories, products and
coproducts coincide and are called \term{biproducts}, written $A \oplus B$.

\begin{theorem}
In a preadditive category $\mathcal{C}$, if a product $A \times B$ and a
coproduct $A + B$ both exist for objects $A, B$, then they are isomorphic.
\end{theorem}

\begin{prooflog}
\proofstep{Write $\iota_1 : A \to A+B$, $\iota_2 : B \to A+B$ for the coproduct
           inclusions, and $\pi_1 : A\times B \to A$, $\pi_2 : A\times B \to B$
           for the product projections.}
           {Standard notation for the universal morphisms.}
\proofstep{The universal property of $A\times B$ gives a unique morphism
           $\phi : A+B \to A\times B$ such that
           $\pi_1 \circ \phi \circ \iota_1 = \operatorname{id}_A$,
           $\pi_2 \circ \phi \circ \iota_2 = \operatorname{id}_B$,
           $\pi_1 \circ \phi \circ \iota_2 = 0$,
           $\pi_2 \circ \phi \circ \iota_1 = 0$.}
           {The cross-terms vanish because the only morphisms involving $0_{a,b}$
            in the preadditive structure allow them to be zero.}
\proofstep{By the bilinearity of composition and the universal property of the
           coproduct, $\phi \circ (\iota_1 \circ \pi_1 + \iota_2 \circ \pi_2) =
           \operatorname{id}_{A\times B}$ and
           $(\iota_1 \circ \pi_1 + \iota_2 \circ \pi_2) \circ \phi =
           \operatorname{id}_{A+B}$.}
           {The sum is well-defined because hom-sets are abelian groups.}
\proofstep{So $\phi$ is an isomorphism $A+B \xrightarrow{\sim} A\times B$.}
           {This isomorphism is the biproduct structure.}
\end{prooflog}

% ─────────────────────────────────────────────────────────────────────────────
\section{Functors Between Algebraic Categories}

Under the algebraic interpretation, functors between one-object categories are
precisely monoid homomorphisms. A functor $F : \mathbf{B}M \to \mathbf{B}N$
maps the single object $\star_M$ to $\star_N$ and acts on morphisms as
$F : M \to N$; functoriality requires $F(m \cdot m') = F(m) \cdot F(m')$ and
$F(e_M) = e_N$. That is precisely a monoid homomorphism.

Similarly, a functor $F : \mathcal{P} \to \mathcal{Q}$ between preorder categories
is an order-preserving map: $a \leq b \Rightarrow F(a) \leq F(b)$.

Natural transformations between functors of one-object categories are
elements of the target monoid satisfying a conjugation condition. A natural
transformation $\alpha : F \Rightarrow G$ between functors
$F, G : \mathbf{B}M \to \mathbf{B}N$ is a single element $n \in N$ such that
for every $m \in M$:
\[
  G(m) \cdot n = n \cdot F(m).
\]
This is a \term{bimodule map condition} or \term{intertwining condition}.

% ─────────────────────────────────────────────────────────────────────────────
\section{Exercises}

\begin{exercisesbox}
\begin{qa}
\qaitem{What algebraic structure corresponds to a one-object category?}
       {A monoid. The single hom-set $\mathcal{C}(\star,\star)$ with composition
        as multiplication and the identity morphism as unit.}

\qaitem{What notation is used for the one-object category associated to a monoid $M$?}
       {$\mathbf{B}M$, the ``delooping'' of $M$.}

\qaitem{What algebraic structure corresponds to a one-object groupoid?}
       {A group. Every morphism has an inverse, so the monoid is a group.}

\qaitem{What is a thin category?}
       {A category where each hom-set has at most one element. Preorders are
        the thin categories.}

\qaitem{What category axioms become when a category is a preorder?}
       {Reflexivity ($a \leq a$) = identity; transitivity = composition. No further axioms.}

\qaitem{What is a groupoid?}
       {A category in which every morphism is an isomorphism.}

\qaitem{How does an equivalence relation embed into category theory?}
       {As a thin groupoid: a thin category where every morphism (arrow $a \to b$
        when $a \sim b$) is invertible.}

\qaitem{What is a preadditive category?}
       {An $\mathbf{Ab}$-enriched category: hom-sets are abelian groups and
        composition is bilinear.}

\qaitem{What is a zero morphism in a preadditive category?}
       {The zero element $0_{a,b} \in \mathcal{C}(a,b)$; it factors through
        any zero object.}

\qaitem{What is a biproduct?}
       {An object $A \oplus B$ that is simultaneously a product and a coproduct
        of $A$ and $B$; exists in additive categories.}

\qaitem{In a preadditive category, when do products equal coproducts?}
       {When both exist, they are isomorphic (the biproduct theorem). The isomorphism
        is constructed via the bilinear structure.}

\qaitem{What is a functor between one-object categories?}
       {A monoid homomorphism. Functoriality is exactly the homomorphism conditions.}

\qaitem{What is a natural transformation between functors of one-object categories?}
       {An element $n$ of the target monoid satisfying $G(m) \cdot n = n \cdot F(m)$
        for all $m$ — an intertwining condition.}

\qaitem{What is a functor between preorder categories?}
       {An order-preserving map: $a \leq b \Rightarrow F(a) \leq F(b)$.}

\qaitem{Why is the algebraic perspective useful?}
       {It shows that classical algebraic structures (monoids, groups, orders) are
        not just analogous to categories but literally special cases, making their
        theorems instances of general categorical theorems.}

\qaitem{What is the category $\mathbf{Gpd}$ of groupoids?}
       {Objects are groupoids; morphisms are functors between them. It contains
        $\mathbf{Grp}$ as a full subcategory via $\mathbf{B}$.}

\qaitem{What does it mean for a category to be ``skeletally small''?}
       {It has only a set of isomorphism classes of objects; equivalently, it is
        equivalent to a small category.}

\qaitem{What is a free category on a directed graph?}
       {Objects are vertices; morphisms are finite paths (sequences of composable edges);
        composition is path concatenation; identity is the empty path at each vertex.}

\qaitem{What universal property does the free category satisfy?}
       {Any function from the graph to the underlying graph of a category $\mathcal{C}$
        extends uniquely to a functor from the free category to $\mathcal{C}$.}

\qaitem{How does the free-forgetful adjunction for categories go?}
       {Free category functor $\mathbf{F} : \mathbf{Grph} \to \mathbf{Cat}$ is left
        adjoint to the underlying-graph functor $U : \mathbf{Cat} \to \mathbf{Grph}$.}

\qaitem{What is a congruence on a category?}
       {An equivalence relation on each hom-set, compatible with composition;
        used to form quotient categories.}

\qaitem{What is the relationship between relations on a monoid and congruences?}
       {A congruence on the one-object category $\mathbf{B}M$ is exactly a congruence
        on the monoid $M$ in the algebraic sense.}
\end{qa}
\end{exercisesbox}

\begin{summarybox}
A category is a \textbf{monoid with multiple objects}: composition is a partially
defined associative operation. One-object categories are exactly monoids; one-object
groupoids are groups; thin categories are preorders. Enriched over $\mathbf{Ab}$,
categories become preadditive; with biproducts and zero object, additive. Functors
between algebraic categories recover algebraic homomorphisms; natural transformations
recover intertwining conditions.
\end{summarybox}

\section{Formal Restatement}

\begin{formalrestatement}
\textbf{Algebraic taxonomy.}
A category $\mathcal{C}$ is:
\begin{itemize}
  \item a \term{monoid} iff $|\operatorname{ob}\mathcal{C}| = 1$;
  \item a \term{group} iff it is a monoid and a groupoid;
  \item a \term{preorder} iff $|\mathcal{C}(a,b)| \leq 1$ for all $a,b$;
  \item a \term{groupoid} iff every morphism is an isomorphism.
\end{itemize}

\textbf{Delooping.} $\mathbf{B} : \mathbf{Mon} \hookrightarrow \mathbf{Cat}$,
$\mathbf{B}M$ has one object $\star$ and $\mathbf{B}M(\star,\star) = M$.

\textbf{Preadditive.} $\mathcal{C}$ is preadditive iff it is $\mathbf{Ab}$-enriched:
each $\mathcal{C}(a,b) \in \mathbf{Ab}$; composition bilinear.

\textbf{Biproduct theorem.} In a preadditive category, if $A\times B$ and $A+B$
both exist, then $A+B \cong A\times B$.

\textbf{Nat transf between $\mathbf{B}M$-functors.} $\alpha : F \Rightarrow G$
for $F, G : \mathbf{B}M \to \mathbf{B}N$ is an element $n \in N$ with
$G(m) \cdot n = n \cdot F(m)$ for all $m \in M$.
\end{formalrestatement}
